%% ============================================================
%% fhe/threshold-custody.tex
%% ============================================================
%%
%% Custody property of the threshold-FHE scheme as deployed in the
%% M-Chain wiring: no member, and no coalition below the threshold,
%% ever holds or can reconstruct the FHE secret key, across key
%% generation and decryption. Also fixes which secret-sharing
%% object the scheme uses and why it is not the one the threshold
%% signature schemes use. Papers cite this; do not redefine.
%% ============================================================

\subsection*{Threshold-FHE custody in the M-Chain wiring}

\begin{definition}[Threshold-FHE committee]
\label{def:tfhe-committee}
A \emph{threshold-FHE committee} is a tuple
$\mathcal{K} = (t, n, \mathsf{pk}, \{\sigma_j\}_{j=1}^{n})$ over the LWE
parameter set $\mathsf{ParameterSetID}$ of
Theorem~\ref{thm:lwe-parameter-security}, where $n$ is the committee size,
$1 \le t \le n$ the threshold, $\mathsf{pk} = (p, a)$ the collective LWE
encryption key, and $\sigma_j$ the share held by the member at position $j$.
Writing $R_q = \mathbb{Z}_q[X]/(X^N+1)$ for the LWE ring, the committee is
\emph{well formed} when there exists $s \in R_q$ and a degree-$(t-1)$ family
$\{f_i\}_{i<N}$ over $\mathbb{Z}_q$ with $f_i(0) = s[i]$ and
$\sigma_j[i] = f_i(j)$ for every $i$ and $j$, and $p = -a s + e$ for an error
$e$ drawn from the parameter set's error distribution.

The value $s$ is the \emph{collective secret}. It is a well-defined element of
$R_q$, determined by any $t$ of the shares; the claim of this section is about
where it does and does not exist as a value in the protocol.
\end{definition}

\begin{definition}[Custody]
\label{def:tfhe-custody}
Let $\Pi$ be a protocol realising Definition~\ref{def:tfhe-committee}, run
over authenticated point-to-point and broadcast channels, and let
$\mathsf{view}_S(\Pi)$ denote the joint view of a coalition $S \subseteq [n]$:
the coalition's own inputs and random tapes, every broadcast message, and
every point-to-point message addressed to a member of $S$. $\Pi$ has
\emph{$t$-custody} when, for every $S$ with $|S| \le t - 1$,
\[
  \mathsf{view}_S(\Pi) \;\approx\; \mathsf{view}_S^{\mathsf{sim}}(\Pi),
\]
where the simulator $\mathsf{view}^{\mathsf{sim}}_S$ is given
$\mathsf{pk}$, the coalition's own contributions, and the plaintexts of the
ciphertexts the committee decrypts during the run, but not $s$. The relation
$\approx$ is computational indistinguishability for the key-generation
transcript and statistical indistinguishability for the decryption transcript.

$t$-custody is strictly stronger than the statement that no \emph{single}
member holds $s$: it also excludes any coalition of $t-1$ members
reconstructing $s$, and it quantifies over the whole run, so a protocol that
avoids materialising $s$ at generation time but reassembles it at decryption
time does not satisfy it.
\end{definition}

\begin{theorem}[Custody of the deployed wiring]
\label{thm:tfhe-mchain-custody}
Let $\Pi_{\mathrm{MPC}}$ be the threshold-FHE lane as deployed: the scheme
$\mathsf{tfhe}$ of \texttt{github.com/luxfi/threshold/protocols/tfhe} ---
dealerless key generation (\texttt{keygen.go}), the partial-decrypt
kernel (\texttt{partial.go}) and the committee layer
(\texttt{committee.go}, \texttt{decrypt.go}) --- driven across nodes by
\texttt{luxfi/mpc/pkg/mpc/tfhe\_session.go} and consumed by the encrypted-policy
gate \texttt{luxfi/mpc/pkg/policy/fhe\_threshold\_decryptor.go}. Then
$\Pi_{\mathrm{MPC}}$ has $t$-custody in the sense of
Definition~\ref{def:tfhe-custody}, under
\begin{enumerate}[nosep]
  \item the decisional RLWE assumption at the deployed parameter set
        (Theorem~\ref{thm:lwe-parameter-security}),
  \item the privacy of Shamir sharing over $\mathbb{Z}_q$ at degree $t-1$, and
  \item the smudging bound of $\mathsf{Sigma}$, namely
        $\Delta \sqrt{t}\,\cdot 6\sigma \le q/16$ with $\Delta = n!$ --- a
        ceiling now conservative by a factor of $\Delta$ relative to what
        recombination needs (Remark~\ref{rem:tfhe-modular-inverse-gap}), and so
        available with room to spare.
\end{enumerate}
\end{theorem}

\begin{proof}
Fix $S$ with $|S| \le t-1$ and let $H = [n] \setminus S$ be the honest
members. The transcript decomposes into the two phases.

\emph{Key generation.} Each member $j$ samples its own contribution $s_j$
locally and publishes $p_j = -a\,s_j + e_j$, where $a$ is the common
reference polynomial. The coalition's view of an honest $j \in H$ therefore
contains: $p_j$, and the single Shamir evaluation $\sigma_{j \to k}[i] =
f_{j,i}(k)$ for each $k \in S$, where $f_{j,i}$ is a fresh degree-$(t-1)$
polynomial with $f_{j,i}(0) = s_j[i]$.

The $|S| \le t-1$ evaluations of a degree-$(t-1)$ polynomial at distinct
non-zero points are uniform on $\mathbb{Z}_q^{|S|}$ and independent of
$f_{j,i}(0)$; this is the privacy half of Shamir sharing, applied
coefficientwise, and it is information-theoretic. The published $p_j$ is an
RLWE sample under $s_j$ and is computationally indistinguishable from uniform
by assumption~(1). Hence the coalition's view of each honest contribution is
simulatable from $\mathsf{pk}$ alone, and so is its view of
$s = \sum_{j} s_j$, of which the coalition contributes only
$\sum_{j \in S} s_j$.

The share the coalition ends up holding is
$\sigma_k[i] = \sum_{j} f_{j,i}(k) = \bigl(\sum_j f_{j,i}\bigr)(k)$, an
evaluation of a degree-$(t-1)$ polynomial whose value at $0$ is $s[i]$. The
coalition holds $|S| \le t-1$ such evaluations, which again determine nothing
about $s[i]$. At no point in this phase does any party receive another
party's contribution, and no party computes $\sum_j s_j$: the sum exists only
as the constant term of a polynomial nobody evaluates at $0$.

\emph{Decryption.} An honest member $j$ answers a decryption request with
$d_j = c_1 \sigma_j + e_j'$, where $e_j'$ is smudging noise at standard
deviation $\sigma$ from $\mathsf{Sigma}(t,n)$. By the noise-flooding
argument of Asharov et al.\ (EUROCRYPT 2012), $d_j$ is statistically close to
a value determined by $c_1$, the plaintext, and the coalition's own shares,
provided the bound in assumption~(3) holds; the coalition therefore learns
nothing about $\sigma_j$ beyond the plaintext it is entitled to. Combination
computes
\[
  \mathsf{pt} \;=\; c_0 \;+\; \sum_{j \in Q} \lambda_j^{Q} d_j
  \qquad (|Q| = t),
\]
with the integer Lagrange coefficients of
Definition~\ref{def:tfhe-recombinable-quorum}. The secret $s$ appears only as
the coefficient of the public $c_1$ inside the masked aggregate; it is never
isolated, so the decryption phase does not reconstruct what the generation
phase declined to form.

Composing the two phases gives $t$-custody.
\end{proof}

\begin{corollary}[No member holds the key]
\label{cor:tfhe-no-member-holds-key}
For $t \ge 2$, no single member of $\Pi_{\mathrm{MPC}}$ holds or can compute
$s$, at any point in the committee's life. This is the $|S| = 1$ case of
Theorem~\ref{thm:tfhe-mchain-custody}.
\end{corollary}

\begin{remark}[Where the property is enforced, not merely claimed]
\label{rem:tfhe-custody-enforcement}
Three code-level facts carry the hypotheses of
Theorem~\ref{thm:tfhe-mchain-custody}, each checked mechanically rather than
by review:
\begin{enumerate}[nosep]
  \item \textbf{No whole key is generated.} \texttt{luxfi/fhe} ships whole-key
        generators ($\mathsf{GenKeyPair}$, $\mathsf{GenSecretKey}$), a
        trusted-dealer split ($\mathsf{Split}$, retained only as a
        known-answer oracle), single-party decryptors and secret-key
        encryptors. None is called from the lane. On the scheme side the
        \texttt{go/ast} gate $\mathsf{TestGate}$
        scans every non-test file of \texttt{protocols/tfhe}; on the wiring
        side $\mathsf{TestNoWholeKeyInThresholdFHESource}$ scans
        \texttt{luxfi/mpc/pkg/mpc/tfhe*.go} and
        \texttt{luxfi/mpc/pkg/policy/fhe\_threshold*.go}. Each has a paired
        negative control showing it flags an injected violation.
  \item \textbf{No whole key crosses a process boundary.} The ceremony's wire
        types ($\mathsf{tfhePublicShare}$, $\mathsf{tfheSubShare}$), the
        scheme's own message types ($\mathsf{Public}$,
        $\mathsf{Point}$, $\mathsf{Share}$,
        $\mathsf{Partial}$) and the persisted member record
        ($\mathsf{Member}$, $\mathsf{Decryption}$) have no secret-key field,
        checked by reflection in
        $\mathsf{TestTransport}$ and
        $\mathsf{TestTFHEWireTypesCarryNoSecretKey}$.
  \item \textbf{Each contribution is erased once spent.} A member calls
        $\mathsf{Party.Zeroize}$ as soon as it holds its share ---
        in $\mathsf{finalizeIfComplete}$, and again in $\mathsf{Stop}$ --- so
        the contribution does not outlive the ceremony.
\end{enumerate}
\end{remark}

\begin{remark}[What changed, and the counterexample it was]
\label{rem:tfhe-prior-construction}
The construction this replaced was not a weaker threshold scheme; it was not a
threshold scheme. Every party stored the whole master key, partial decryption
returned a MAC tag, and combination ignored the partials and ran single-party
decryption. It therefore satisfied correctness --- $t$ parties recovered the
plaintext --- while failing Definition~\ref{def:tfhe-custody} for every
$t \ge 2$ by the trivial distinguisher: a coalition of size $1$ holds $s$, so
its view is not simulatable without $s$.

It is worth stating as a counterexample because it is the shape of failure that
correctness testing cannot see. A round-trip test over that construction
passed. It was removed from \texttt{github.com/luxfi/threshold} for that
reason, and the scheme described here now occupies the same import path.
\end{remark}

\begin{remark}[Why this scheme sits with the other threshold schemes]
\label{rem:tfhe-with-the-others}
$\mathsf{tfhe}$ lives in \texttt{luxfi/threshold} beside threshold ECDSA
(\texttt{protocols/cmp}, \texttt{protocols/doerner}), threshold Schnorr
(\texttt{protocols/frost}), threshold BLS (\texttt{protocols/bls}) and the
post-quantum threshold signatures (\texttt{protocols/corona},
\texttt{protocols/mldsa}, \texttt{protocols/pulsar}). What it shares with them
is the property being proved and the structure that delivers it: a committee of
$n$ members at threshold $t$; a generation round in which each member
contributes only its own secret and the group secret is never formed; a use
round in which each member emits a masked partial result and the group secret
is never reassembled. Definition~\ref{def:tfhe-custody} is the same definition
those schemes are held to, with signing replaced by decryption. Keeping them
apart invited the failure in
Remark~\ref{rem:tfhe-prior-construction}: a scheme kept outside the module
where that law is enforced was not measured against it.

What it does not share is the arithmetic, and that boundary is stated in
Remark~\ref{rem:tfhe-secret-sharing} rather than papered over.
\end{remark}

\begin{remark}[Which secret sharing, and why there are three]
\label{rem:tfhe-secret-sharing}
Three distinct linear secret-sharing objects appear in this stack and are not
interchangeable. Recording which is which is load-bearing, because the
correctness of threshold decryption depends on the third and on nothing else.
\begin{enumerate}[nosep]
  \item \textbf{Curve-scalar LSS} (\texttt{luxfi/threshold/protocols/lss}).
        The field is an elliptic-curve group order, the Feldman commitments are
        curve points, reconstruction uses modular inverses. This is the sharing
        the threshold signature schemes use.
  \item \textbf{Coefficientwise $\mathbb{Z}_q$ sharing}
        (\texttt{protocols/tfhe}, $\mathsf{deal}$). A share is a
        polynomial in $R_q$ shared coefficientwise over $\mathbb{Z}_q$, where
        $q$ is the LWE modulus of the parameter set --- about $2^{27}$ at
        PN10QP27. Reconstruction must avoid modular inverses entirely: the
        combine scales each contribution by the \emph{integer} numerator
        $\Lambda_j = \Delta \lambda_j$, $\Delta = n!$, precisely so that the
        smudging noise $e_j'$ is multiplied by a bounded integer rather than by
        an arbitrary field element. That is what keeps
        $|\sum_j \Lambda_j e_j'|$ inside the decoding budget and makes
        assumption~(3) of Theorem~\ref{thm:tfhe-mchain-custody} available.
  \item \textbf{A 2048-bit MODP commitment sharing}
        (\texttt{luxfi/fhe/pkg/keycommit}). Shamir with exponentiation-based
        Feldman commitments over RFC~3526 Group~14, applied to the SHA-256
        commitment of a key's components. It shares a hash, so its shares
        cannot decrypt; it is a dealer-equivocation audit facility and is
        named for that.
\end{enumerate}
Objects (1) and (2) are the same idea over different fields with different
reconstruction, so they are not merged: sharing the FHE secret in a curve
scalar field would share a different secret, and clearing denominators with
modular inverses would destroy Theorem~\ref{thm:tfhe-mchain-custody}'s noise
bound and with it correctness. Object (3) is not a threshold scheme at all and
therefore does not belong beside either.
\end{remark}

\begin{definition}[Recombinable quorum]
\label{def:tfhe-recombinable-quorum}
For a quorum $S \subseteq [n]$ the Lagrange coefficient at zero of member
$i \in S$ is
\[
  \lambda_i^S \;=\; \prod_{j \in S,\, j \neq i} \frac{x_j}{x_j - x_i},
  \qquad x_k = k .
\]
$S$ is \emph{recombinable} when $\lambda_i^S \in \mathbb{Z}$ for every
$i \in S$. Not every subset is: for $S = \{1,3\}$ one has
$\lambda^S = (-3/2,\, 1/2)$.
\end{definition}

\begin{theorem}[Correctness on recombinable quorums, and refusal otherwise]
\label{thm:tfhe-recombinable-correctness}
Let $\mathsf{ct} = (c_0, c_1)$ encrypt $m$ under a well-formed committee
(Definition~\ref{def:tfhe-committee}) and let each $i \in S$ contribute
$d_i = c_1 \sigma_i + e_i'$. If $S$ is recombinable then
\[
  c_0 + \sum_{i \in S} \lambda_i^S d_i
  \;=\; m + e + \sum_{i \in S} \lambda_i^S e_i' \pmod q,
\]
with $|\lambda_i^S| \le \binom{n-1}{t-1}$, and the plaintext decodes correctly
whenever that noise stays inside the decoding margin. If $S$ is not
recombinable, $\mathsf{combine}$ returns an error and no plaintext.
\end{theorem}

\begin{proof}
The displayed identity is linearity of the sharing:
$\sum_{i \in S} \lambda_i^S \sigma_i = s$ by Lagrange interpolation at zero,
so the $c_1$ terms collapse to $c_1 s$ and $c_0 + c_1 s = m + e$. The
coefficient bound is the standard bound on Lagrange coefficients over
$\{1,\dots,n\}$. The second clause is by construction: the implementation
computes each $\lambda_i^S$ as an exact integer quotient and returns an error
when any remainder is non-zero, before touching the ciphertext.
\end{proof}

\begin{remark}[Gap found and closed: recombination by modular inverse]
\label{rem:tfhe-modular-inverse-gap}
The previous recombination did clear the denominator, and that was a
correctness defect rather than a noise-budget question. It followed the
Bendlin--Damg\aa rd presentation: recombine the scaled secret with integer
numerators $\Lambda_i^S = \Delta \lambda_i^S$, $\Delta = n!$, obtaining
$\Delta c_0 + \sum_i \Lambda_i^S d_i = \Delta(m+e) + \sum_i \Lambda_i^S e_i'$,
and then recover $m$ by multiplying by $\Delta^{-1} \bmod q$.

That last step is sound on the message term and destroys the noise term. The
noise $\sum_i \Lambda_i^S e_i'$ is a multiple of $\Delta$ only when
$\Delta \mid \Lambda_i^S$ for every $i$, which is exactly the condition
$\lambda_i^S \in \mathbb{Z}$ of Definition~\ref{def:tfhe-recombinable-quorum}.
When it fails, multiplication by $\Delta^{-1} \bmod q$ maps a small noise value
to an arbitrary field element, the rounding margin is gone, and the recovered
bit is wrong --- deterministically, with no error raised.

The effect was measured before it was fixed, at both deployed parameter sets.
Of the quorums of threshold size, those that recombine and those that decrypted
correctly coincided exactly: at $2$-of-$3$, $2$ of $3$ quorums; at $3$-of-$4$,
$3$ of $4$; at $3$-of-$5$, $5$ of $10$; at $4$-of-$7$, $13$ of $35$; at
$5$-of-$9$, $45$ of $126$. The non-recombinable quorums returned the wrong
plaintext for one or both plaintext polarities. Widening the modulus from
$q \approx 2^{27}$ (PN10QP27) to $q \approx 2^{54}$ (PN11QP54) left the pattern
unchanged, which is what distinguishes this from a noise-budget failure.

It had gone unnoticed because the end-to-end tests recombined the quorum of the
$t$ largest positions, $\{n-t+1, \dots, n\}$, which is recombinable for every
$(t,n)$ --- a single favourable ordering standing in for the property.

The fix is to not clear denominators at all: accept exactly the recombinable
quorums, recombine with the integer $\lambda_i^S$ directly per
Theorem~\ref{thm:tfhe-recombinable-correctness}, and refuse the rest. This
also strictly improves the noise, since $|\lambda_i^S| \le \binom{n-1}{t-1}$
replaces $|\Lambda_i^S| \le \Delta\binom{n-1}{t-1}$. Tests:
$\mathsf{TestQuorate}$,
$\mathsf{TestQuorum}$ (for every quorum:
correct value or an error, never a different value), and
$\mathsf{TestRefusal}$.
\end{remark}

\begin{remark}[The access structure that results, and its price]
\label{rem:tfhe-access-structure}
Restricting recombination to recombinable quorums does not weaken
Theorem~\ref{thm:tfhe-mchain-custody}: custody is a statement about what
$t-1$ members learn, and it is unchanged. What is restricted is
\emph{availability}. The realised access structure is not every $t$-subset but
the recombinable ones, so a committee whose only $t$ responsive members form a
non-recombinable set refuses the decryption instead of answering it.

Two facts keep that price small, and both are checked. The quorum of the $t$
largest positions is recombinable for every $(t,n)$ in range
($\mathsf{TestLargest}$, exhaustive
for $n \le 24$), so a committee with all members responsive always has a usable
quorum; and both consumers search the contributions they hold for a
recombinable subset rather than fixing one in advance
($\mathsf{Decrypt}$ via $\mathsf{choose}$, and the policy dispatcher via
$\mathsf{Quorum}$, which keeps a round open for one more member instead of
refusing at the threshold count).

Restoring the full $t$-of-$n$ structure means removing the division rather than
guarding it: either scale the plaintext by $\Delta$ at encryption time, so that
recovering $\Delta m$ needs no correction, at the cost of $\log_2 \Delta$ bits
of decoding margin; or replace Shamir with a monotone-span-program sharing
whose reconstruction coefficients are $0/1$, at the cost of
$\binom{n}{t-1}$-sized shares. Both are parameter and encoding changes with
their own noise accounting, and neither is what the implementation does today.
\end{remark}

\begin{remark}[Gap found and closed: ciphertext binding]
\label{rem:tfhe-ciphertext-binding}
Theorem~\ref{thm:tfhe-mchain-custody} is a secrecy statement and holds without
any binding between a partial decryption and its ciphertext. Correctness of a
\emph{networked} combination does not.

$\mathsf{combine}$ checks that the partials it is given agree on the LWE
modulus and the committee size and carry distinct member positions. It cannot
check which ciphertext a partial was computed for, because a partial carries no
such information. A member --- or a relay --- that answers a request about
$\mathsf{ct}$ with a partial decryption of some other $\mathsf{ct}'$ under the
same committee key therefore contributes a well-formed-looking value that
combines into a plaintext bearing no relation to either ciphertext. In the
policy gate the recovered plaintext is a single verdict bit, so this is a
one-bit corruption of an authorisation decision.

The committee layer closes it by binding. $\mathsf{Decryption}$ carries
$\mathsf{Digest} = H(\text{\texttt{LUX/THRESHOLD/TFHE/CIPHERTEXT/v1}} \,\|\,
\mathsf{ct})$, computed inside $\mathsf{Member.Decrypt}$ from the same bytes it
decrypts, so a producer cannot make the two disagree. $\mathsf{Decrypt}$
refuses any contribution whose digest is not the digest of the ciphertext under
decryption, and the policy dispatcher ($\mathsf{Committee.collect}$)
additionally declines to count such a contribution towards the quorum, so one
mismatching member does not deny a round that honest members can still carry.
Tests: $\mathsf{TestBinding}$,
$\mathsf{TestCommittee\_ReplayedPartialDoesNotFillQuorum}$,
$\mathsf{TestCommittee\_ReplayedPartialRefusesWhenItIsTheQuorum}$.
\end{remark}

\begin{remark}[Residual: no proof of correct partial decryption]
\label{rem:tfhe-active-corruption}
Custody and robustness are different properties and only the first is proved
here. Theorem~\ref{thm:tfhe-mchain-custody} bounds what a coalition
\emph{learns}; it says nothing about what a coalition can make the committee
\emph{output}. A member that submits $d_j' = d_j + \delta$ for a $\delta$ of
its choosing, correctly bound to the right ciphertext, shifts the combined
plaintext by $\Lambda_j \delta$, and nothing in the scheme detects it: a
partial decryption carries no proof that it was computed under the member's own
share. Closing this needs a proof of correct partial decryption --- the
CDS/NIZK argument that accompanies the Asharov et al.\ scheme in the malicious
model --- which this scheme does not implement. The behaviour is pinned by
$\mathsf{TestTampered}$, which fails if an
integrity check is ever added, forcing this remark to be revised with it.

Authentication does not substitute for it. The ceremony and decryption messages
are Ed25519-signed on the session wire ($\mathsf{session.sendMsg}$ /
$\mathsf{verifyInboundSignature}$), which establishes \emph{which} member sent a
partial decryption and nothing about whether that member computed it honestly.
Attestation of the host (\texttt{luxfi/mpc/cc/attest}) is a third, separate
property and likewise does not imply correct computation. The deployed integrity
assumption is therefore: custody against any $t-1$ passive corruptions, and
correctness of the recovered plaintext only when every member in the responding
quorum is honest.
\end{remark}

\begin{remark}[Common reference polynomial]
\label{rem:tfhe-crs}
The reference polynomial $a$ is derived by
$\mathsf{tfheCRSSeed} = H(\text{\texttt{LUX/MPC/TFHE/DKG/CRS/v1}} \,\|\,
\mathsf{walletID} \,\|\, \langle \text{committee in canonical order}\rangle)$
and expanded with a keyed PRNG, so every node derives the identical $a$ from
local information and no agreement round is needed; $a$ is public and its only
requirement is agreement. Modelled as a random-oracle output, $a$ is uniform on
$R_q$ as assumption~(1) requires. Whoever chooses the wallet label influences
$a$ only by selecting among random-oracle outputs, which is the standard
hash-derived-CRS model; a beacon-supplied seed would remove even that.
Determinism is load-bearing and is tested
($\mathsf{TestTFHECRSSeed\_KnownAnswer}$,
$\mathsf{TestTFHECRSSeed\_PermutationInvariant}$,
$\mathsf{TestTFHECeremony\_ReferencePolynomialFollowsTheSeed}$): two nodes
disagreeing on $a$ produce a committee that agrees on no public key.
\end{remark}

\begin{remark}[Scope: encryption and decryption, not homomorphic compute]
\label{rem:tfhe-lane-scope}
The dealerless lane generates the collective \emph{encryption} key and its
$t$-of-$n$ shares. It does not generate the FHEW blind-rotation key, and
dealerless generation of that key --- a separate multiparty-FHEW protocol ---
is a documented residual. Consequently the committee can trustlessly encrypt
and trustlessly threshold-decrypt, which is what
Theorem~\ref{thm:tfhe-mchain-custody} covers and what the wallet path uses, but
homomorphic \emph{evaluation} on ciphertexts under the collective key still
requires an evaluation key produced outside the lane. Any deployment that
expects the encrypted policy verdict of
\texttt{luxfi/mpc/pkg/policy} to be computed homomorphically under the
committee's collective key depends on that residual; the decryptor itself is
correct for any ciphertext under the collective key, however produced.
\end{remark}
